%%%%%%%%%%%%%%%%%%%%%%%%%%%%%%%%%%%%%%%%%%%%%%%%%%%%%%%
% % Lines starting with % are comments, which are ignored.
% % This is a handy way of indicating the date and version of
% % your document, to wit:
% %
% % LaTeX sample file
% % Modified March, 2002
% %
%%%%%%%%%%%%%%%%%%%%%%%%%%%%%%%%%%%%%%%%%%%%%%%%%%%%%%%
\DocumentMetadata{
  % PDF standards: ONE key whose value is a comma-separated list --
  % UA-2 (accessibility) plus A-4f (archival; the archival companion of
  % UA-2). Do NOT use separate pdfstandard=... keys: with ordinary
  % key-value parsing only the last one takes effect.
  pdfstandard={UA-2,A-4f},
  lang=en-US,
  tagging=on,
  tagging-setup={
    math/setup=mathml-SE
  }
}

\documentclass{ltx-talk}
% lecture-macros.tex -- shared package setup for lecture-template.tex
% (ltx-talk slides / handouts).
%
% NOTE ON PATHS: every template in this repository is stored one level up
% from where it is compiled (see README.md). All relative \input paths
% below are resolved against the folder the document is compiled in,
% which sits one directory above this file.

% --- Tables, bookmarks, strikethrough ------------------------------------
% booktabs: professional table rules (\toprule, \midrule, \bottomrule)
% bookmark: PDF outline entries (PDF bookmarks)
% ulem[normalem]: \sout strikethrough that survives math mode
\usepackage{booktabs}
\usepackage{bookmark}
\usepackage[normalem]{ulem}
\usepackage{multicol}

% --- Math ------------------------------------------------------------
% amsmath: equation environments and numbering
% mathrsfs: \mathfrak  |  stmaryrd: extra symbols
\usepackage{amsmath}
\usepackage{mathrsfs}
\usepackage{stmaryrd}

% --- Layout -------------------------------------------------------------
% float: [H] float placement
% fontspec: Unicode/OpenType fonts (compile with XeLaTeX or LuaLaTeX)
\usepackage{float}
\usepackage{fontspec}

% --- Colors ---------------------------------------------------------------
% dvipsnames: extended named color palette (the Accessible* colors below
% are built on it);
% table: \rowcolors / \cellcolor for shading tables.
% NOTE: xcolor may only be loaded ONCE per document -- a second
% \usepackage{xcolor} with different options is a fatal "Option clash"
% error, so ALL of its options are listed here.
\usepackage[dvipsnames,table]{xcolor}

% --- Multimedia -----------------------------------------------------------
% animate: embeds audio/video (and legacy Flash) media in a frame.
\usepackage{animate}

% --- Bibliography ----------------------------------------------------------
% biblatex (biber backend): citations for lecture notes.
\usepackage[
  backend=biber,
  style=alphabetic,
  url=true,
  doi=true,
  isbn=true,
  backref=true
]{biblatex}

% --- Colored boxes ----------------------------------------------------------
% talkthemetcolorbox: the colored-box theme for ltx-talk slides
% (shadow/rounded/shaded blocks, showtitle, titlepage styles).
\usepackage[
  shadow,
  rounded,
  shaded,
  blocks,
  showtitle,
  titlepage
]{talkthemetcolorbox}

% --- "theindex" helper environment ------------------------------------------
% Provides a lightweight theindex environment (an itemize) and the
% \@idxitem/\subitem/\subsubitem/\indexspace macros that LaTeX's standard
% theindex machinery expects, so hand-written index pages indent correctly.
\makeatletter
\newenvironment{theindex}
  {\begin{itemize}}
  {\end{itemize}}
\providecommand\@idxitem{\par\hangindent 40\p@}
\providecommand\subitem{\@idxitem \hspace*{20\p@}}
\providecommand\subsubitem{\@idxitem \hspace*{30\p@}}
\providecommand\indexspace{\par \vskip 10\p@ \@plus5\p@ \@minus3\p@\relax}
\makeatother

% --- Index -----------------------------------------------------------------
% imakeidx: builds the printed index.
% The redefinition below appends "|hyperpage" to every index entry so that,
% in the PDF, each index entry is a link to its page.
\usepackage{imakeidx}
\makeatletter
\let\imki@orig@wrindexentrysplit\imki@wrindexentrysplit
\renewcommand\imki@wrindexentrysplit[3]{%
  \imki@orig@wrindexentrysplit{#1}{#2|hyperpage}{#3}%
}

\makeatletter
% --- Section titles: capture + one bookmark per logical section -------------
% \section is redefined so that:
%   1. the section title is captured into the plain-TeX register
%      \mytalksectitletoks (reprinted verbatim on the auto divider frame,
%      see the cmd/section/after hook below -- plain-TeX storage is
%      catcode-safe, unlike expanding the title in a hook), and
%   2. exactly ONE depth-1 PDF bookmark is created per \section call
%      (hyperref's own bookmark is suppressed via \Hy@writebookmark).
\DeclareCommandCopy{\OriginalSectionForBookmarks}{\section}
\newtoks\mytalksectitletoks

\RenewDocumentCommand{\section}{s o m}{%
  \begingroup
    \def\Hy@writebookmark##1##2##3##4##5{}%
    \global\mytalksectitletoks={#3}%
    \IfBooleanTF{#1}
      {\phantomsection\IfNoValueTF{#2}{\OriginalSectionForBookmarks*{#3}}{\OriginalSectionForBookmarks*[#2]{#3}}}%
      {\IfNoValueTF{#2}{\OriginalSectionForBookmarks{#3}}{\OriginalSectionForBookmarks[#2]{#3}}}%
  \endgroup
  \IfBooleanTF{#1}
    {\IfNoValueTF{#2}{\pdfbookmark[1]{#3}{section-star.\thepage.\number\inputlineno}}{\pdfbookmark[1]{#2}{section-star.\thepage.\number\inputlineno}}}%
    {\IfNoValueTF{#2}{\pdfbookmark[1]{Section \thesection: #3}{section.\thesection}}{\pdfbookmark[1]{Section \thesection: #2}{section.\thesection}}}%
}
\makeatother

% --- One PDF bookmark per logical frame (NOT per overlay) ---------------------
% ltx-talk writes one PDF page per overlay, so a naive frametitle hook would
% register a new bookmark for every overlay of the same frame. The
% \g_slideutils_frame_bookmark_int register remembers the last frame that
% got a bookmark, so subsequent overlays of the same frame are skipped.
% This is also what keeps BeamerPresenter's cumulative annotations working:
% one logical frame must stay one logical unit in the PDF outline.
\ExplSyntaxOn
\int_new:N \g_slideutils_frame_bookmark_int

\AddToHook{cmd/frametitle/after}{%
  \tl_if_empty:NF \g__talk_frame_title_tl
    {
      \ifnum\value{frame}=\g_slideutils_frame_bookmark_int\relax
      \else
        \int_gset:Nn \g_slideutils_frame_bookmark_int { \value{frame} }%
        \pdfbookmark[2]{\tl_use:N \g__talk_frame_title_tl}
          {frame.\int_use:N \g_slideutils_frame_bookmark_int}
      \fi
    }
}
\ExplSyntaxOff
\makeatother

% --- Printed index + theme for the colored boxes -------------------------------
% \makeindex[...]: two-column index titled "Index", listed in the TOC.
% The two EditInstance calls style the standard header/footer instances
% of talkthemetcolorbox (cyan background; footer order title/author/
% date/framenumber).
\makeindex[columns=2, title=Index, intoc]\EditInstance{header}{std}{
  background-color = cyan,
  color            = black,
  height           = 0.75cm
}
\EditInstance{footer}{std}{
  background-color = cyan,
  element-order = {title, author, date, framenumber},
}

% --- Auto section divider frame ---------------------------------------------------
% After each \section, insert a centered, large, structure-colored frame
% that prints the section title (from \mytalksectitletoks).
\AddToHook{cmd/section/after}{%
  \begin{frame}
    \centering
    \color{structure}
    \LARGE
    \the\mytalksectitletoks
  \end{frame}%
}

% --- Tagging sockets (tagged PDF structure) ---------------------------------------
% ltx-talk exposes tagging "sockets" where the PDF structure can attach
% section titles; register one and leave it unassigned (none) so the
% class's own tagging keeps control.
\ExplSyntaxOn
\AtBeginDocument{
  \NewTaggingSocketPlug { talk / sec / title } { none } { }
  \AssignTaggingSocketPlug { talk / sec / title } { none }
}
\ExplSyntaxOff

% --- Accessibility roles ----------------------------------------------------------------
% Map \frametitle to the H1 heading role in the tagged PDF, so screen
% readers treat slide titles as headings.
\tagpdfsetup{role / new-tag = frametitle / H1}

\makeatletter
% --- Page label = frame number -------------------------------------------------------
% Redefine what feeds the PDF page label to track the frame counter,
% not the physical page counter. Replace \theframenumber below with
% whatever internal macro/property ltx-talk actually uses to produce
% the "framenumber" footer element — confirm via ltx-talk-frame.dtx.
\renewcommand\thepage{\arabic{frame}}
\makeatother

\addbibresource{refs.bib}

% --- "Problem Time!" slides ------------------------------------------------------------------
% \ProblemSlide{problem}{solution} emits two frames: the problem, then
% its solution after a \pause, each numbered by the problem counter.
\newcounter{problemcounter}
\setcounter{problemcounter}{0}
\newcommand{\ProblemSlide}[2]{
  \stepcounter{problemcounter}
  \begin{frame}
    \frametitle{Problem Time!}
    Problem \theproblemcounter: #1
  \end{frame}

  \begin{frame}
    \frametitle{Problem Time!}
    Solution to Problem \theproblemcounter: \pause

    #2
  \end{frame}
}
% \highlight[<colour>]{<stuff>}
% from https://tex.stackexchange.com/questions/33401/a-version-of-colorbox-that-works-inside-math-environments
\newcommand{\highlight}[2][yellow]{\mathchoice%
  {\colorbox{#1}{$\displaystyle#2$}}%
  {\colorbox{#1}{$\textstyle#2$}}%
  {\colorbox{#1}{$\scriptstyle#2$}}%
  {\colorbox{#1}{$\scriptscriptstyle#2$}}}%

% --- Fonts ------------------------------------------------------------------------------
% Lato (text) + LeteSansMath (math): open, accessible typefaces.
\setmainfont{Lato}[
  Ligatures=TeX,
  Numbers={Lining,Monospaced}
]
\setsansfont{Lato}[
  Ligatures=TeX,
  Numbers={Lining,Monospaced}
]
\setmathfont{LeteSansMath.otf}\usepackage{emoji}

% --- Accessible colors ---------------------------------------------------------------------
% WCAG-contrast-safe variants of red/green for text on white:
% red!80!black and green!40!black both clear the 4.5:1 requirement.
% AltYellow: soft background for highlight boxes.
\colorlet{AccessibleRed}{red!80!black}
\colorlet{AccessibleGreen}{green!40!black}
\definecolor{AltYellow}{RGB}{255,245,170}

% --- Links ---------------------------------------------------------------------------------
% hyperref + link colors: every link uses AccessibleGreen (7.2:1 on
% white), so internal, citation, and URL links all pass WCAG 1.4.3.
\usepackage{hyperref}
\hypersetup{
  colorlinks=true,
  linkcolor=AccessibleGreen,
  urlcolor=AccessibleGreen,
  citecolor=AccessibleGreen,
}

% --- Shared course macros ---------------------------------------------------------------------
% Root macros.tex: math shorthands used across all course documents
% (\RR, \NN, \ZZ, \QQ, \grad, \curl, \tr, \der, \pder, ...).
\newcommand{\RR}{\mathbb{R}}
\newcommand{\NN}{\mathbb{N}}
\newcommand{\ZZ}{\mathbb{Z}}
\newcommand{\QQ}{\mathbb{Q}}
\newcommand{\vv}{\vec{v}}
\newcommand{\ww}{\vec{w}}
\newcommand{\xx}{\vec{x}}
\newcommand{\yy}{\vec{y}}
\newcommand{\zz}{\vec{z}}
\newcommand{\uu}{\vec{u}}
\newcommand{\basis}[1]{\vec{e_{#1}}}
\newcommand{\id}{\mathrm{id}}
\newcommand{\CC}{\mathbb{C}}
\renewcommand{\div}{\operatorname{div}}
\newcommand{\grad}{\operatorname{grad}}
\newcommand{\curl}{\operatorname{curl}}
\newcommand{\tr}{\operatorname{tr}}
\newcommand{\der}[2]{\dfrac{d #1}{d #2}}
\newcommand{\pder}[3][]{\dfrac{\partial^{#1} #2}{\partial #3}}


% % Title and author(s)
%%%%%%%%%%%%%%%%%%%%%%%%%%%%%%%%%%%%%%%%%%%%%%%%%%%%%%%
\title{Lecture 4: Limits III}
\author{Garrett Credi}
\date{09/17/2026}
%%%%%%%%%%%%%%%%%%%%%%%%%%%%%%%%%%%%%%%%%%%%%%%%%%%%%%%
% %
% % The next command allows your in import encapsulated
% % postscript files, .epsf or .eps files, which
% % contain vector graphic image data.
% %
%%%%%%%%%%%%%%%%%%%%%%%%%%%%%%%%%%%%%%%%%%%%%%%%%%%%%%%
% %   The first thanks indicates your affiliation
% %
% %  Just the name here.
% %
% % Your mailing address goes at the end.
% %
% % \thanks is also how you indicate grant support
% %
%%%%%%%%%%%%%%%%%%%%%%%%%%%%%%%%%%%%%%%%%%%%%%%%%%%%%%%

%%%%%%%%%%%%%%%%%%%%%%%%%%%%%%%%%%%%%%%%%%%%%%%%%%%%%%%
% $ The next command ensures the margins are 1 inch on
% % all sides.
% %
%%%%%%%%%%%%%%%%%%%%%%%%%%%%%%%%%%%%%%%%%%%%%%%%%%%%%%%
\begin{document}
\begin{frame}
  \maketitle
\end{frame}

\begin{frame}
  \frametitle{Quiz Time!}
  Let's get into formation for quiz time.
\end{frame}

\begin{frame}
  \frametitle{Last Time}

  We spent much of last lecture developing our ability to compute limits via: \pause

  \begin{enumerate}
    \item Recognizing when functions were \textit{continuous} so that we could reduce the computation of limits to simpler evaluations \pause
    \item Developing limit properties which allowed us to build up the computation of complicated limits via reduction to simpler ones 
  \end{enumerate}

  We also developed the notion of the \textit{one sided limits} - restricted kinds of limits where we need only approach from a singular side.
\end{frame}

\begin{frame}
  \frame{Today's Goals}

  We'll apply both our new limit computation skills as well as our knowledge of one-sided limits to work with \textit{limits at or involving infinity}.

  We will also uncover our last tool to help us compute limits: \textit{the Squeeze Theorem.}
\end{frame}

\section{Limits at/involving Infinity}

\begin{frame}
  \frametitle{Prelude I}

  Let's first examine one of the problems from the previous worksheet: \pause

  We had two functions whose graphs looked like: \pause
  \vspace{2in}

  And we wanted to compute $\lim_{x \rightarrow 1} g(f(x)) \pause = \lim_{x \rightarrow 1} g((x-1)^2) \pause = \lim_{u \rightarrow 0^+} g(u) = 3$
\end{frame}

\begin{frame}
  \frametitle{Prelude II} 

  Key to that problem was noticing that: \pause
  \begin{equation}
    x \rightarrow 1 \pause \quad \iff \quad (x-1) \rightarrow 0^+ \pause
  \end{equation}

  Which allowed us to do a ``substitution" that made the limit easier to compute. \pause

  This will be our strategy for computing limits at $\pm \infty$! \pause
  Since $\pm \infty$ is not an actual number we can compute with in a consistent/coherent way, \pause
  we'll try and replace $x \rightarrow \pm \infty$ with a limit as \textit{some other quantity} approaches \textit{some actual number}.
\end{frame}

\begin{frame}
  \frametitle{Limits at Infinity I}

  Before that, let's actually look at the literal definition to see why we want to avoid it: \pause
  \begin{block}{Limits at Infinity}\index{limits!at infinity}
    For a function $f(x)$, we say that $\lim_{x \rightarrow \infty} f(x) = L$ \pause
    if we can make $f(x) \approx L$ (up to some error $\epsilon > 0$) just by forcing $x > N$ for some real number $N$. \pause

    Similarly, we say that $\lim_{x \rightarrow -\infty} f(x) = L$ \pause
    if we can make $f(x) \approx L$ (up to some error $\epsilon > 0$) just by forcing $x < M$ for some real number $M$. \pause
  \end{block}

  Do we want to prove things in this class like we would need to if we actually used this definition? \pause
  \textcolor{red}{NO!}
\end{frame}

\begin{frame}
  \frametitle{Limits at Infinity II}
  Instead, let's see if we can relate it to something we could instead \textit{compute}
  using our tools from last lecture. \pause

  \begin{align}
    x > N &\iff 0 > x > N \pause \\
          &\iff 0 < \dfrac{1}{x} < \dfrac{1}{N} \pause \\
          &\iff 0 < \dfrac{1}{x} - 0 < \dfrac{1}{N} \pause
  \end{align}

  In other words: \pause
  \begin{alertblock}{Limits at Infinity in terms of Finite Limits}\index{limits!at infinity!in terms of finite limits}
    Taking the limit as $x \rightarrow + \infty$ is \textit{exactly equivalent} to taking the limit as $\dfrac{1}{x} \rightarrow 0^+$ \pause

    Similarly $x \rightarrow - \infty$ is \textit{exactly equivalent} to $\dfrac{1}{x} \rightarrow 0^{-}$!
  \end{alertblock}
\end{frame}

\begin{frame}
  \frametitle{Limits at Infinity III: Example}

  For example: \pause
  \begin{align}
    \lim_{x \rightarrow \infty} \dfrac{3x^2 + 2x + 1}{2x^2 + 6x + 131} &= \pause \lim_{\frac{1}{x} \rightarrow 0^+} \dfrac{3x^2 + 2x + 1}{2x^2 + 6x + 131} \pause \\
                                                                       &= \lim_{u \rightarrow 0^+} \pause \dfrac{3(\frac{1}{u})^2 + 2(\frac{1}{u}) + 1}{2(\frac{1}{u})^2 + 6(\frac{1}{u}) + 131} \pause \\
                                                                       &= \lim_{u \rightarrow 0^+} \dfrac{u^2}{u^2} \dfrac{3 \frac{1}{u^2} + 2\frac{1}{u} + 1}{2 \frac{1}{u^2} + 6\frac{1}{u} + 131} \pause \\
                                                                       &= \lim_{u \rightarrow 0^+} \dfrac{3 + 2u + u^2}{2 + 6u + 131u^2} \pause \\
                                                                       &= \lim_{u \rightarrow 0} \dfrac{3 + 2u + u^2}{2 + 6u + 131u^2} \pause = \dfrac{3 + 2\cdot 0 + 0^2}{2 + 6\cdot 0 + 131 \cdot 0^2} = \pause \dfrac{3}{2} \pause
  \end{align}
\end{frame}
\begin{frame}
  \frametitle{Limits at Infinity IV: Example 1+2}
  So we took a scary looking limit \pause $\lim_{x \rightarrow \infty} \dfrac{3x^2 + 2x + 1}{2x^2 + 6x + 131}$ \pause
  and reduced it to simply evaluating some rational function \pause $f(u) = \dfrac{3 + 2u + u^2}{2 + 6u + 131u^2}$ at $u = 0$! \pause

  What about for just a polynomial? \pause
  \begin{align}
    \lim_{x \rightarrow - \infty} (x^3 + 2x + 1) &= \pause \lim_{u \rightarrow 0^{-}} ((\dfrac{1}{u})^3 + 2 \dfrac{1}{u} + 1) \pause \\
                                                 &= \lim_{u \rightarrow 0^{-}} \dfrac{1}{u^3} (1 + 2u^2 + u^3) \pause  \\
    ``&=" (\lim_{u \rightarrow 0^{-}} \dfrac{1}{u^3}) \cdot (\lim_{u \rightarrow 0^{-}} (1 + 2u^2 + u^3)) \pause \\
                                                 &= \lim_{u \rightarrow 0^{-}} \dfrac{1}{u^3} \pause
  \end{align}

  Where we unfortunately encounter a problem \pause since $\dfrac{1}{u^3}$ has an \textit{asymptote}
  at $u = 0$ \pause and so cannot be replaced by a function continuous at $0$. \pause What do we do?
\end{frame}

\begin{frame}
  \frametitle{Limits at Infinity V: Example 2 and Divergence to Infinity}
  Let's look at the graph! \pause

  \begin{block}{Limits Diverging to Infinity}\index{limits!diverging to infinity}
    For a function $f(x)$ and a point $a$, we say that $\lim_{x \rightarrow a} f(x) = +\infty$ (read as the limit of $f(x)$ as $x$ approaches $a$ diverges to positive infinity) \pause
    if, for any lower bound $L$ \pause we can force $f(x) > L$ \pause just by forcing $x \approx a$ (up to some error $\delta > 0$). \pause

    Similarly, we say that $\lim_{x \rightarrow a} f(x) = - \infty$ (... diverges to minus infinity) \pause
    if, for any upper bound $U$ \pause we can force $f(x) < U$ \pause just by forcing $x \approx a$ (up to some error $\delta > 0$). \pause
  \end{block}

  So, from the graph of $\dfrac{1}{u^3}$ as $u \rightarrow 0^{-}$ we see that \pause
  $\lim_{u \rightarrow 0^{-}} \dfrac{1}{u^3}$ \textit{diverges} \pause to minus infinity.
\end{frame}

\section{Squeeze Theorem}

\begin{frame}
  \frametitle{Squeeze Theorem I}

  Let's consider one final limit: \pause
  \begin{align}
    \lim_{x \rightarrow \infty} \dfrac{\sin(x)}{x} &= \pause \lim_{u \rightarrow 0^{+}} \dfrac{\sin(\frac{1}{u})}{\frac{1}{u}} \pause \\
                                                   &= \lim_{u \rightarrow 0^{+}} u \sin(\frac{1}{u}) \pause
  \end{align}

  And now we \textit{REALLY} get stuck: \pause as $u \rightarrow 0$, $\sin(\frac{1}{u})$
  oscillates infinitely rapidly and escapes all ability to be pinned down to a value we could use to simplify the limit. \pause
  Noooooo! \pause
  
  ... \pause ... \pause but are we really stuck?
\end{frame}

\begin{frame}
  \frametitle{Squeeze Theorem II}

  How much does $\sin(\frac{1}{u})$ really affect the value of the limit? \pause
  Since we know that $\sin$ as a function never gets large at all, \pause
  intuitively it shouldn't be able to outpace $u$ shrinking towards $0$. \pause

  Formally, since $-1 \leq \sin(x) \leq 1$ (the range of $\sin(x)$ is $[-1,1]$) \pause
  we have both that: \pause
  \begin{align}
    \lim_{u \rightarrow 0^{+}} u \sin(\frac{1}{u}) &\leq \lim_{u \rightarrow 0^{+}} u \cdot \pause 1 \pause \\
                                                   &= \lim_{u \rightarrow 0^{+}} u \pause = \lim_{u \rightarrow 0} u \pause \\
                                                   &= 0 \pause
  \end{align}
  and that: \pause
  \begin{align}
    \lim_{u \rightarrow 0^{+}} u \sin(\frac{1}{u}) & \geq \lim_{u \rightarrow 0^{+}} u \cdot \pause -1 \pause \\
                                                   &= - \lim_{u \rightarrow 0^{+}} u \pause = - \lim_{u \rightarrow 0} u \pause \\
                                                   &= - 0 = 0 \pause
  \end{align}
\end{frame}

\begin{frame}
  \frametitle{Squeeze Theorem III}
  So if \pause
  \begin{equation}
  0 \leq \lim_{u \rightarrow 0^{+}} u \sin(\frac{1}{u}) \leq 0 \pause
  \end{equation}

  we must have that $\lim_{u \rightarrow 0^{+}} u \sin(\frac{1}{u}) = 0$! \pause
  
  This style of argument comes from the following result: \pause
  \begin{alertblock}{The Squeeze Theorem}\index{Squeeze Theorem}\index{limits!Squeeze Theorem}
    Suppose we have a function $f(x)$ and a point $a$ so that: \pause
    \begin{enumerate}
      \item we can find other functions $g(x), h(x)$ so that \pause (around $x = a$) $g(x) \leq f(x) \leq h(x)$ \pause
      \item and that $\lim_{x \rightarrow a} g(x) = \lim_{x \rightarrow a} h(x)$ \pause
    \end{enumerate}

    Then we must have that $\lim_{x \rightarrow a} f(x)$ exists \pause and is equal to the
    common value of $\lim_{x \rightarrow a} g(x) = \lim_{x \rightarrow a} h(x)$.
  \end{alertblock}
\end{frame}

\begin{frame}
  \frametitle{Tools for Algebraically Computing Limits}

  Let's recap our tools: \pause
  \begin{enumerate}
    \item Evaluation for continuous functions: \pause if $f(x)$ is continuous at $x = a$ then \pause
      \begin{equation}
        f(a) \quad \pause = \quad \lim_{x \rightarrow a^{+}} f(x) = \lim_{x \rightarrow a} f(x) = \lim_{x \rightarrow a^{-}} f(x)
      \end{equation} \pause
    \item Reduction to continuous functions: \pause maybe the point at which you are taking the limit towards isn't in the domain \pause
      but you can replace it with one and use tool \#1
    \item Limit Properties: \pause You can rewrite your limit as a sum, product, quotient, or composite of other limits 
      and reduce it to computing those \pause
    \item Inversion for Limits at infinity: \pause $x \rightarrow \infty \iff \frac{1}{x} \rightarrow 0^{+}$, $x \rightarrow -\infty \iff \frac{1}{x} \rightarrow 0^{-}$ \pause
    \item Squeeze Theorem: \pause maybe you have a limit which involves a term that ``shouldn't matter" intuitively \pause
      and you can use the range of that term to put bounds that squeeze the value of the limit.
  \end{enumerate}
\end{frame}

\begin{frame}
  \frametitle{Worksheet Time}

  Let's get into groups and work on the worksheet!
\end{frame}

\begin{frame}
  \frametitle{End of Class}

  That's all folks!
\end{frame}
\end{document}
